% Non-SNEA Algonquian Resources — External Resources
% Draftable now. Full details in research cards at .issues/research-cards/

\section{Non-SNEA Algonquian Resources}
\label{sec:external}

This project draws on a range of Algonquian resources beyond the SNEA languages
directly represented in the corpus. These are organized by language family group
below. Full details for each resource are documented in the research cards at
\texttt{.issues/research-cards/}.

\subsection{Proto-Algonquian Reconstructions}

\begin{itemize}
  \item \textbf{Bloomfield (1925, 1946):} Foundational PA reconstruction.
  \item \textbf{Aubin (1975):} \emph{A Proto-Algonquian Dictionary.} Mercury
    Series, Canadian Ethnology Service Paper No. 39. \textasciitilde2,300
    reconstructions. Analyzed Narragansett as substantially accurate relative
    to PA.
  \item \textbf{Hewson (1993):} \emph{A Computer-Generated Dictionary of
    Proto-Algonquian.} 4,066 entries. Online at
    \url{https://protoalgonquian.atlas-ling.ca}. Typos corrected 2015--2017.
  \item \textbf{Pentland (1979):} \emph{Algonquian Historical Phonology.}
    Comprehensive PA-to-daughter phonological developments.
\end{itemize}

\subsection{Eastern Algonquian (Closest Genetic Relatives)}

Eastern Algonquian is the closest genetic relative of the SNEA languages.
These resources provide the best comparative data for resolving SNEA
phonological questions.

\begin{itemize}
  \item \textbf{Mi'kmaq Talking Dictionary} (\url{https://mikmaqonline.org})
    --- 6,000+ words, each recorded by 3+ speakers. Preserves PA cluster
    distinctions collapsed in SNEA colonial sources.
  \item \textbf{Passamaquoddy-Maliseet Dictionary} (Francis \& Leavitt,
    \url{https://pmportal.org}) --- 18,000+ words with video archive. 50+
    years of collaboration among speakers, educators, and linguists.
  \item \textbf{Western Abenaki Dictionary} (Day 1994, 2 vols, 538 pp) ---
    Documents the language as spoken in the last half of the 20th century.
    Abenaki is the closest living relative of the SNEA languages.
  \item \textbf{Nulhegan Abenaki Dictionary} (2025,
    \url{https://abenakionline.com}) --- First publicly shareable English-Abenaki
    dictionary, released with National Archives grant support.
  \item \textbf{Lenape Talking Dictionary} (\url{https://www.talk-lenape.org})
    --- 17,000+ words with audio, stories, usage examples. Lenape is
    transitional between SNEA and Central Algonquian.
  \item \textbf{A Lenâpé-English Dictionary} (Brinton, from Moravian MS,
    \url{https://archive.org}) --- 19th-century dictionary based on Zeisberger,
    Heckewelder, and other Moravian sources.
\end{itemize}

\subsection{Central Algonquian (Most Conservative Phonologically)}

Central Algonquian languages are the most conservative phonologically and
provide the best evidence for PA reflex verification.

\subsubsection*{Cree-Innu Group}

\begin{itemize}
  \item \textbf{Eastern James Bay Cree Dictionary}
    (\url{https://dictionary.eastcree.org}) --- Northern and Southern dialects,
    English-Cree-French.
  \item \textbf{itwêwina Plains Cree Dictionary}
    (\url{https://itwewina.altlab.app}) --- Searchable Plains Cree (Y-dialect).
  \item \textbf{Online Cree Dictionary} (\url{https://www.creedictionary.com})
    --- Aggregates 4 dictionaries: Alberta Elders', Earle Waugh, Wolvengrey,
    Maskwacis.
  \item \textbf{A Dictionary of the Cree Language} (Watkins 1865/1938,
    \url{https://www.creeliteracy.org}) --- Historical dictionary for
    diachronic comparison.
\end{itemize}

\subsubsection*{Ojibwe-Potawatomi Group}

\begin{itemize}
  \item \textbf{Ojibwe People's Dictionary}
    (\url{https://ojibwe.lib.umn.edu}) --- 7,000+ entries with audio.
    Southwestern Ojibwe (N-dialect).
  \item \textbf{Nishnaabemwin Online Dictionary}
    (\url{https://dictionary.nishnaabemwin.atlas-ling.ca}) --- 12,000+ words.
    Odawa dialects of Lake Huron.
  \item \textbf{A Dictionary of the Ojibway Language} (Baraga 1853/1992,
    \url{https://archive.org}) --- Historical dictionary by Bishop Frederic
    Baraga.
  \item \textbf{Potawatomi Language Dictionary}
    (\url{https://www.potawatomidictionary.com}) --- Citizen Potawatomi Nation.
\end{itemize}

\subsubsection*{Sauk-Fox-Kickapoo Group}

\begin{itemize}
  \item \textbf{A Concise Dictionary of the Sauk Language}
    (\url{https://www.sacandfoxnation-nsn.gov}) --- Based on fieldwork
    1990s--2000s.
  \item \textbf{Kickapoo Vocabulary} (Voorhis 1988) --- Algonquian and
    Iroquoian Linguistics, Memoir 6. 205 pp.
\end{itemize}

\subsubsection*{Miami-Illinois and Other Central}

\begin{itemize}
  \item \textbf{Myaamia ILDA Dictionary}
    (\url{https://aacimotaatiiyankwi.org}) --- Miami-Illinois online dictionary
    with mobile app. Myaamia Center, Miami University.
  \item \textbf{The Menomini Language} (Bloomfield 1962,
    \url{https://archive.org}) --- 515 pp. The most phonologically conservative
    Algonquian language. Preserves PA vowel length and *h preaspiration most
    faithfully.
  \item \textbf{English-Shawnee Dictionary} (Emerson) --- Practical dictionary.
\end{itemize}

\subsection{Plains Algonquian (Highly Divergent)}

Plains Algonquian languages are the most divergent from PA and provide the
outer bounds for PA reflex verification.

\begin{itemize}
  \item \textbf{Blackfoot Online Dictionary}
    (\url{https://dictionary.blackfoot.atlas-ling.ca}) --- Part of the
    Algonquian Dictionaries Project.
  \item \textbf{English-Cheyenne Dictionary} (Petter 1915,
    \url{https://cheyennelanguage.org}) --- 1,126 pp with handwritten
    annotations. Cheyenne shows extreme PA consonant shifts.
  \item \textbf{Dictionary of the Arapaho Language} (Salzmann 1983, Cowell
    et al.\ 2012, \url{https://www.colorado.edu/project/arapaho}) --- 243 pp
    (4th ed.). Arapaho shows extreme PA vowel reduction.
\end{itemize}

\subsection{Comparative SNEA Dialectology}

\begin{itemize}
  \item \textbf{Costa (2007):} ``The Dialectology of Southern New England
    Algonquian.'' Definitive modern dialect classification of SNEA.
  \item \textbf{Goddard \& Bragdon (1988):} \emph{Native Writings in
    Massachusett.} Comparative phonology/orthography.
  \item \textbf{Bragdon (2009):} Updated linguistic analysis of Massachusett.
  \item \textbf{Cunningham (2024):} Nanticoke phonological analysis ---
    directly parallel methodology.
  \item \textbf{Masthay (2001):} Mahican language word list.
  \item \textbf{Goddard (1978--2010):} Eastern Algonquian surveys, Munsee
    Delaware phonology, Mahican phoneme inventory.
  \item \textbf{Siebert (1975):} Methods for reconstructing from early
    documentary sources (Powhatan).
\end{itemize}
